\documentclass[conference]{IEEEtran}
\IEEEoverridecommandlockouts
% The preceding line is only needed to identify funding in the first footnote. If that is unneeded, please comment it out.
\usepackage{cite}
\usepackage{amsmath,amssymb,amsfonts}
\usepackage{algorithmic}
\usepackage{algorithm}
\usepackage{graphicx}
\usepackage{textcomp}
\usepackage{xcolor}
\def\BibTeX{{\rm B\kern-.05em{\sc i\kern-.025em b}\kern-.08em
    T\kern-.1667em\lower.7ex\hbox{E}\kern-.125emX}}

\begin{document}

\title{Exploring Parameter-Efficient Adaptation for Time Series Foundation Models}

\author{\IEEEauthorblockN{Emanuele Romito}
\IEEEauthorblockA{
\textit{Politecnico di Torino}\\
Turin, Italy \\
s339170@studenti.polito.it}
}

\maketitle

\begin{abstract}
\end{abstract}

\section{Introduction}

Time series forecasting has evolved from classical statistical methods such as ARIMA \cite{box2015time} to deep learning approaches. The introduction of Transformer-based architectures marked a significant paradigm shift, with models like Informer \cite{zhou2021informer}, Autoformer \cite{wu2021autoformer}, and PatchTST \cite{nie2023patchtst} demonstrating the efficacy of attention mechanisms for long-sequence modeling. This progression has culminated in the emergence of general-purpose Time Series Foundation Models (TSFMs). Unlike specific models trained from scratch, TSFMs leverage large-scale pre-training on diverse datasets to learn universal temporal patterns, enabling zero-shot transfer capabilities across unseen domains.

Among these emerging foundation models, Chronos-2 \cite{ansari2025chronos2univariateuniversalforecasting} represents a distinct paradigm, adopting a language-modeling approach by quantizing time series values into discrete tokens. A key innovation is its "Group Attention" mechanism, which enables the model to process multiple multivariate series simultaneously, thereby capturing cross-series dependencies. While Chronos-2 exhibits impressive zero-shot performance, the transition to foundation models introduces new challenges in deployment and adaptation.

While TSFMs face various hurdles in real-world application, this research focuses on a specific opportunity for expansion. To address the issue of non-stationarity where the standard "pretrain-then-inference" assumption fails, we incorporate PETSA \cite{medeiros2025accurateparameterefficienttesttimeadaptation} for parameter-efficient test-time adaptation.

\section{Problem Statement}

We define the time series forecasting problem under the setting of Parameter-Efficient Test-Time Adaptation (PETSA). The goal is to predict future values of a target time series.

\subsection{PETSA}
\textbf{Task:} The objective of this task is to adapt a pre-trained TSFM to a batch of target time series concurrently during inference. This allows for instance-specific adaptation, where each time series in the batch receives its own unique set of adapted parameters to handle its specific non-stationarity or domain shift.

\textbf{Input:}
\begin{itemize}
    \item A batch of $B$ target time series contexts $\mathcal{X} = \{\mathbf{x}^{(i)}_{1:T}\}_{i=1}^{B}$.
    \item A pre-trained foundation model parameterized by $\theta$.
\end{itemize}

\textbf{Output:}
\begin{itemize}
    \item A set of $B$ instance-specific adapted parameter sets $\{\theta'^{(i)}\}_{i=1}^{B}$, where $\theta'^{(i)} = \theta + \Delta\theta^{(i)}$. Each $\Delta\theta^{(i)}$ represents the update specific to the $i$-th time series (e.g., via batched LoRA adapters), obtained by minimizing the forecasting error on $\mathbf{x}^{(i)}_{1:T}$.
    \item The forecasts $\{\mathbf{x}^{(i)}_{T+1:T+H}\}_{i=1}^{B}$ generated using their respective adapted models.
\end{itemize}

\section{Methodology}

To facilitate the exploration of the proposed extension (PETSA), we first undertook the replication of the Chronos-2 training pipeline. This foundational step was necessary to establish a robust baseline and ensuring that our modifications were applied to a verified codebase, consistent with the reference architecture.

\subsection{Chronos-2 Training Pipeline Overview}
We implemented a modular training framework designed to replicate the Chronos-2 setup. The system consists of the following key components:

\begin{enumerate}
    \item \textbf{Synthetic Data Generation:} Generation of diverse time series using KernelSynth (Gaussian Processes with mixed kernels) to capture universal temporal dynamics.
    \item \textbf{Dataset Building:} Construction of a \texttt{StreamingChronosDataset} that dynamically samples from synthetic and real-world sources.
    \item \textbf{Model Definition:} Implementation of the Chronos-2 architecture based on the T5 encoder backbone, augmented with Group Attention.
    \item \textbf{Training:} Execution of the optimization loop minimizing Weighted Quantile Loss (WQL) with checkpointing and metric logging.
    \item \textbf{Benchmarking:} Validation of the baseline model using a standardized suite before extending the pipeline with the PETSA module.
\end{enumerate}

\subsection{PETSA Implementation}
We extended the Chronos-2 foundation model by implementing a bespoke adaptation pipeline. Our implementation diverges from standard fine-tuning by utilizing a "Batched LoRA" approach, enabling efficient, instance-specific adaptation for multiple distinct time series in a single forward pass.

\subsubsection{Architecture Description}
The extension is built upon three primary modules:

\begin{itemize}
    \item \textbf{BatchedLoRALayer:} A custom linear layer wrapper that modifies the standard Low-Rank Adaptation (LoRA) formulation. While standard LoRA learns a single pair of matrices $A$ and $B$ for the entire dataset, our layer instantiates batched parameters $\mathbf{A} \in \mathbb{R}^{B \times r \times d_{in}}$ and $\mathbf{B} \in \mathbb{R}^{B \times d_{out} \times r}$. This allows the layer to compute unique weight updates $\Delta W^{(i)} = \alpha (\mathbf{B}^{(i)} \mathbf{A}^{(i)})$ for each sample $i$ in the batch simultaneously via vectorized batch matrix multiplication, maintaining high training throughput.
    
    \item \textbf{ChronosPETSAWrapper:} This module serves as the functional interface to the frozen backbone. During initialization, it inspects the Chronos-2 encoder and replaces the query, key, and value projections in the \texttt{TimeSelfAttention} mechanism with \texttt{BatchedLoRALayer} instances. It encapsulates the adaptation logic, managing the lifecycle of the batched parameters and the gradient-based optimization loop.
    
    \item \textbf{ChronosPETSAPipeline:} A high-level pipeline that orchestrates the data flow. It extends the standard inference pipeline to support adaptation by generating self-supervised tasks from the test context (masking the recent history) and invoking the wrapper's optimization routine before generating the final forecast.
\end{itemize}

\subsubsection{Pseudocode}
The core logic of the Batched PETSA process is outlined in Algorithm \ref{alg:petsa}.

\begin{algorithm}
\caption{Batched Test-Time Adaptation Strategy}
\begin{algorithmic}
\REQUIRE Batch contexts $\mathbf{X} \in \mathbb{R}^{B \times L}$, Horizon $H$, Steps $K$
\ENSURE Forecasts $\mathbf{Y} \in \mathbb{R}^{B \times H}$
\STATE Initialize batched params $\mathbf{A}, \mathbf{B}$ for all LoRA layers
\STATE Freeze backbone $\theta$, set $\mathbf{A}, \mathbf{B}$ differentiable
\FOR{step $k=1$ to $K$}
    \STATE Split $\mathbf{X}$ into $\mathbf{X}_{in}, \mathbf{X}_{tgt}$ (mask last $M$ tokens)
    \STATE $\hat{\mathbf{X}} \leftarrow \text{Model}(\mathbf{X}_{in}; \theta, \mathbf{A}, \mathbf{B})$
    \STATE $\mathcal{L} \leftarrow \text{WQL}(\hat{\mathbf{X}}, \mathbf{X}_{tgt})$
    \STATE Compute gradients $\nabla_{\mathbf{A}, \mathbf{B}} \mathcal{L}$
    \STATE Update $\mathbf{A}, \mathbf{B}$ using AdamW
\ENDFOR
\STATE $\mathbf{Y} \leftarrow \text{Model}(\mathbf{X}; \theta, \mathbf{A}, \mathbf{B})$ (predict $H$ steps)
\RETURN $\mathbf{Y}$
\end{algorithmic}
\label{alg:petsa}
\end{algorithm}

\section*{Appendix: Model Configuration}
\begin{table}[htbp]
\caption{Chronos-2 (Base) Model Configuration}
\begin{center}
\begin{tabular}{|l|c|}
\hline
\textbf{Parameter} & \textbf{Value} \\
\hline
Parameters & 120M \\
Encoder & T5 \\
Layers & 12 \\
Attention Heads & 12 \\
Hidden Dimension ($d_{model}$) & 768 \\
Feed-Forward Dimension ($d_{ff}$) & 3072 \\
Context Length & 8192 tokens \\
Prediction Horizon & 1024 tokens \\
Patch Size & 16 \\
Dropout & 0.1 \\
\hline
\end{tabular}
\label{tab:model_config}
\end{center}
\end{table}

\section*{Appendix: Training Datasets}
\label{section:dataset-appendix}
The following real-world datasets were used in the training mixture:
\begin{itemize}
    \item \textbf{Monash Repository:} Australian Electricity, Electricity Hourly/Weekly, KDD Cup 2018, Temperature Rain, London Smart Meters.
    \item \textbf{M4 Competition:} Daily, Hourly, Monthly, Weekly.
    \item \textbf{Transportation:} Mexico City Bikes, Taxi (1h, 30min), Uber TLC (Daily, Hourly).
    \item \textbf{Energy & Weather:} Solar (1h), USHCN Daily, Weatherbench (Daily, Hourly, Weekly), Wind Farms (Daily, Hourly).
    \item \textbf{Other:} Wiki Daily 100k.
    \item \textbf{GiftEval (Zero-Shot):} Alibaba Cluster Trace, Azure VM Traces, Borg Cluster Data, Q-Traffic, Buildings 900k.
\end{itemize}

\bibliographystyle{IEEEtran}
\bibliography{references}

\end{document}
